% =============================================================================
%  Viscosity Measurement Laboratory - Student Handout
%  A. Bulent Koc, Clemson University
%  DOI 10.5281/zenodo.22218165
%
%  Compile:  pdflatex viscosity-lab-handout.tex   (twice, for the page count)
%  Paper:    A4, 2.0 cm margins, page numbers bottom centre as "n of N"
% =============================================================================
\documentclass[11pt,a4paper]{article}

\usepackage[a4paper,top=2.1cm,bottom=2.2cm,left=2.1cm,right=2.1cm,
            headheight=14pt,footskip=1.1cm]{geometry}
\usepackage[T1]{fontenc}
\usepackage[utf8]{inputenc}
\usepackage{lmodern}
\usepackage{amsmath,amssymb}
\usepackage{array}
\usepackage{booktabs}
\usepackage{colortbl}
\usepackage{enumitem}
\usepackage{xcolor}
\usepackage{fancyhdr}
\usepackage{lastpage}
\usepackage{tikz}
\usepackage{microtype}
\usepackage{needspace}
\usepackage[colorlinks=true,linkcolor=black,urlcolor=inkblue,citecolor=black]{hyperref}

% ---- colours ---------------------------------------------------------------
\definecolor{inkblue}{HTML}{1A3F7A}   % anything the student writes in
\definecolor{ink2}{HTML}{3D444B}
\definecolor{ink3}{HTML}{767D84}
\definecolor{rulegrey}{HTML}{B9BCC0}
\definecolor{notebg}{HTML}{F4F2EC}
\definecolor{cautionbg}{HTML}{FBF3EC}
\definecolor{cautionrule}{HTML}{A8642A}
\definecolor{gridminor}{HTML}{DDE1E5}
\definecolor{gridmajor}{HTML}{B9BFC5}

% ---- running head and foot -------------------------------------------------
\pagestyle{fancy}
\fancyhf{}
\fancyhead[L]{\footnotesize\color{ink3} Viscosity Measurement Laboratory}
\fancyhead[R]{\footnotesize\color{ink3} Name \rule[-0.15em]{4.2cm}{0.4pt}}
\fancyfoot[C]{\footnotesize\color{ink2} Page \thepage\ of \pageref{LastPage}}
\renewcommand{\headrulewidth}{0.4pt}
\renewcommand{\footrulewidth}{0pt}
\renewcommand{\headrule}{\hbox to\headwidth{\color{rulegrey}\leaders\hrule height \headrulewidth\hfill}}

\fancypagestyle{firstpage}{%
  \fancyhf{}%
  \fancyfoot[C]{\footnotesize\color{ink2} Page \thepage\ of \pageref{LastPage}}%
  \renewcommand{\headrulewidth}{0pt}%
}

% ---- section styling -------------------------------------------------------
\usepackage{titlesec}
\titleformat{\section}{\normalfont\large\bfseries}{\thesection}{0.6em}{}
\titleformat{\subsection}{\normalfont\normalsize\bfseries\itshape}{\thesubsection}{0.6em}{}
\titlespacing*{\section}{0pt}{1.6ex plus .3ex}{0.7ex}
\titlespacing*{\subsection}{0pt}{1.2ex plus .2ex}{0.5ex}

% ---- fill-in blanks --------------------------------------------------------
% \bl   short blank   \bll  medium   \blll long
\newcommand{\bl}{\rule[-0.35em]{1.6cm}{0.4pt}}
\newcommand{\bll}{\rule[-0.35em]{2.6cm}{0.4pt}}
\newcommand{\blll}{\rule[-0.35em]{4.6cm}{0.4pt}}
\newcommand{\bls}{\rule[-0.35em]{1.0cm}{0.4pt}}
% a full-width ruled answer line
\newcommand{\answerline}{\par\vspace{0.9em}\noindent\rule{\linewidth}{0.4pt}}
\newcommand{\answerlines}[1]{\foreach \i in {1,...,#1}{\answerline}}

% ---- note boxes ------------------------------------------------------------
\usepackage[most]{tcolorbox}
\newtcolorbox{labnote}[1][]{
  enhanced, breakable, colback=notebg, colframe=ink3,
  boxrule=0pt, leftrule=2pt, arc=0pt, outer arc=0pt,
  left=8pt, right=8pt, top=6pt, bottom=6pt,
  fontupper=\small, #1}
\newtcolorbox{labcaution}[1][]{
  enhanced, breakable, colback=cautionbg, colframe=cautionrule,
  boxrule=0pt, leftrule=2pt, arc=0pt, outer arc=0pt,
  left=8pt, right=8pt, top=6pt, bottom=6pt,
  fontupper=\small, #1}
\newcommand{\noteh}[1]{{\footnotesize\bfseries\sffamily\color{ink3}\MakeUppercase{#1}}\par\vspace{2pt}}
\newcommand{\cautionh}[1]{{\footnotesize\bfseries\sffamily\color{cautionrule}\MakeUppercase{#1}}\par\vspace{2pt}}

% ---- tables ----------------------------------------------------------------
\renewcommand{\arraystretch}{1.55}
\newcolumntype{C}[1]{>{\centering\arraybackslash}p{#1}}
\newcolumntype{L}[1]{>{\raggedright\arraybackslash}p{#1}}
% interior rules light, outer rules black
\newcommand{\rowrule}{\arrayrulecolor{rulegrey}\hline\arrayrulecolor{black}}

% ---- graph paper -----------------------------------------------------------
\newcommand{\graphpaper}[2]{%
  \begin{tikzpicture}
    \draw[step=0.25cm,gridminor,very thin] (0,0) grid (#1,#2);
    \draw[step=1.25cm,gridmajor,thin]      (0,0) grid (#1,#2);
    \draw[black,thick] (0,0) rectangle (#1,#2);
  \end{tikzpicture}}

\setlist[itemize]{leftmargin=1.3em,itemsep=1pt,topsep=2pt}
\setlist[enumerate]{leftmargin=1.6em,itemsep=10pt,topsep=3pt}

\begin{document}
\thispagestyle{firstpage}

% =============================================================================
\begin{center}
  {\LARGE\bfseries Viscosity Measurement:\\[2pt]
   Falling Ball, Saybolt, and Rotary Viscometers\par}
  \vspace{5pt}
  {\large\itshape Student Handout\par}
  \vspace{3pt}
  {\normalsize Fluid Power Systems and Internal Combustion Engines\par}
  \vspace{4pt}
  {\footnotesize\sffamily\color{ink3}
   AGRICULTURAL MECHANIZATION AND BUSINESS \quad$\cdot$\quad CLEMSON UNIVERSITY\par}
\end{center}
\vspace{6pt}

\noindent
\begin{tabular}{@{}p{0.11\linewidth}p{0.36\linewidth}p{0.09\linewidth}p{0.36\linewidth}@{}}
Name  & \rule{\linewidth}{0.4pt} & Date  & \rule{\linewidth}{0.4pt} \\
Group & \rule{\linewidth}{0.4pt} & Bench & \rule{\linewidth}{0.4pt} \\
\end{tabular}

\vspace{4pt}
\noindent\rule{\linewidth}{1pt}
\vspace{2pt}

% =============================================================================
\section{Objectives}
\begin{itemize}
  \item To become familiar with viscosity measurement methods.
  \item To measure viscosity of hydraulic fluids with falling ball viscometer, Saybolt
        viscometer and rotary viscometer.
  \item To convert one set of viscosity units to another set of units.
  \item To analyze the effect of temperature on viscosity of hydraulic fluid.
\end{itemize}

\section{Core Principle}
Viscosity is a measure of a fluid's resistance to flow. None of the three instruments on this
bench measures that resistance directly. Each one measures something else, how fast a sphere
sinks, how long a fixed volume takes to drain, how hard a rotor must be driven, and then infers
viscosity from a model. The model is where the physics lives, and the model can fail while the
instrument still gives you a perfectly repeatable number.

Two quantities carry the name viscosity and they are not interchangeable. \emph{Dynamic} (or
absolute) viscosity $\mu$ measures resistance to shear and is reported in Pa$\cdot$s or
centipoise. \emph{Kinematic} viscosity $\nu$ is dynamic viscosity divided by density and is
reported in m$^2$/s or centistokes. Density is the only bridge between them. An ISO hydraulic
grade such as ISO~46 refers to kinematic viscosity in centistokes at 40\,\textdegree C; a
Stokes' law result is dynamic viscosity in Pa$\cdot$s. Quoting one where the other is expected
is the most common error in this laboratory.

\begin{labnote}
\noteh{What you should be able to say at the end}
Which of your three measurements you trust most, and why. That answer depends on the Reynolds
number you compute in Section~5.6, not on which instrument looked most expensive.
\end{labnote}

\section{Safety}
\begin{itemize}
  \item Fluids heated to 50 to 60\,\textdegree C will burn skin. Use gloves and tongs on the
        Saybolt cup and the rotary beaker.
  \item Wipe every spill at once. Oil on a hard floor is a fall hazard.
  \item Do not heat any fluid above the temperature stated in the procedure.
  \item Recover the steel sphere with the magnet. Never invert the falling ball tube.
  \item Used fluid goes in the marked waste container, never down the sink.
  \item Report a cracked tube or a leaking cup to the instructor before proceeding.
\end{itemize}

\needspace{5\baselineskip}
\section{Constants and Conversions}
Record these once and use them throughout.

\smallskip
\noindent{\small
$g = 9.81$\,m/s$^2$ \quad$\cdot$\quad 1\,in $=$ 2.54\,cm \quad$\cdot$\quad
1\,m $=$ 100\,cm \quad$\cdot$\quad 1000\,g $=$ 1\,kg \quad$\cdot$\quad
1\,m$^3$ $=$ 1000\,L $=$ $10^6$\,cm$^3$\par
\smallskip
1\,poise $=$ 0.1\,Pa$\cdot$s $=$ 100\,cP \quad$\cdot$\quad
1\,cP $=$ 0.001\,Pa$\cdot$s \quad$\cdot$\quad
1\,Pa$\cdot$s $=$ 1000\,cP $=$ 1000\,mPa$\cdot$s \quad$\cdot$\quad
cSt $=$ cP\,/\,$\rho$\,[g/cm$^3$]\par}

% =============================================================================
\needspace{6\baselineskip}
\section{Part I: Falling Ball Viscometer}
One of the oldest ways of measuring the viscosity of a fluid is to measure how fast a sphere
falls through it. A sphere released in a viscous fluid accelerates for a short distance and then
settles at a constant \emph{terminal velocity}, where the drag force exactly balances the
sphere's weight less its buoyancy. If the sphere falls quickly the fluid has low viscosity and
little resistance to flow; if it falls slowly the fluid has high viscosity. Stokes' law relates
that settling velocity to viscosity, rearranged below to solve for viscosity.

\begin{equation}
  \mu = \frac{2\,\Delta\rho\,g\,r^{2}}{9\,v}
\end{equation}
\noindent{\small
\begin{tabular}{@{\hspace{1.5em}}l@{\quad}l@{}}
$\mu$          & absolute (dynamic) viscosity, Pa$\cdot$s \\[-6pt]
$\Delta\rho$   & density of sphere minus density of fluid, kg/m$^3$ \\[-6pt]
$g$            & acceleration due to gravity, 9.81\,m/s$^2$ \\[-6pt]
$r$            & radius of sphere, m \\[-6pt]
$v$            & velocity of sphere, m/s \\
\end{tabular}\par}

\begin{labcaution}
\cautionh{Release the sphere above the START mark}
Equation~(1) applies only to a sphere already at terminal velocity. If you release the sphere at
the START mark it is still accelerating as it crosses it, your timed velocity is too low, and
your viscosity comes out too high. Release it at the top of the tube and let it settle before it
reaches the mark. The bench sphere needs about 0.25\,m of travel to reach terminal velocity in a
motor oil and about 0.5\,m in an ATF, so the START mark must be at least that far below the
release point.
\end{labcaution}

\begin{labnote}
\noteh{The whole station in one multiplication}
Once the sphere, the marks and the fluid density are fixed, Equation~(1) collapses to
$\mu = K\,t_\mathrm{avg}$ with $K = 2\,\Delta\rho\,g\,r^2/(9L)$. For the bench sphere and marks, $K$
is about 0.68\,Pa$\cdot$s per second for both bench fluids. Compute $K$ once, then every fall time
becomes a viscosity in one step. Section~5.6 reduces the same way:
$\mathit{Re} = \rho_f L (2r)/(K\,t_\mathrm{avg}^2)$, about $14/t_\mathrm{avg}^2$ on this bench.
\end{labnote}

\subsection{Steel sphere}
Measure the diameter of the sphere with a micrometer and its mass on the balance. Compute the
volume from Equation~(2) and the density from Equation~(3).

\begin{equation}
  V = \tfrac{4}{3}\pi r^{3}
\end{equation}
\begin{equation}
  \rho = \frac{m}{V}
\end{equation}

\noindent Mass of sphere: \bl\ g $=$ \bl\ kg \par
\smallskip
\noindent Diameter $D$: \bl\ in $=$ \bl\ cm $=$ \bl\ m \hfill Radius $r$: \bl\ m \par
\smallskip
\noindent Volume $V$: \bll\ m$^3$ \hfill Sphere density $\rho_s$: \bll\ kg/m$^3$ \par

\subsection{Fluid density}
Fill a beaker with 50\,mL of each fluid and weigh it. Record the temperature with a thermometer.
Divide the measured mass in grams by the volume in millilitres to obtain the density, then carry
it through to kg/m$^3$.

\begin{center}
\needspace{6\baselineskip}
{\small
\begin{tabular}{L{0.30\linewidth} C{0.28\linewidth} C{0.28\linewidth}}
\toprule
\textbf{Quantity} & \textbf{Fluid 1} & \textbf{Fluid 2} \\
\midrule
Fluid name                  & & \\ \rowrule
Temperature, \textdegree C  & & \\ \rowrule
Beaker volume, mL           & & \\ \rowrule
Fluid mass, g               & & \\ \rowrule
$\rho$, g/mL                & & \\ \rowrule
$\rho$, g/L                 & & \\ \rowrule
$\rho$, kg/m$^3$            & & \\
\bottomrule
\end{tabular}\par
\vspace{2pt}
{\footnotesize\itshape Table 1. Fluid identification and density.}}
\end{center}

\subsection{Fall time and velocity}
Measure the distance between the two marked points on the viscometer. Hold the sphere with the
magnet, release it from the top of the tube, and time its travel between the two marks. Repeat
four times and average. Four trials are taken because a single stopwatch reading carries a
reaction-time error of roughly a tenth of a second, which on a two-second fall is a five percent
error on its own.

\smallskip
\noindent Distance between marks $L$: \bl\ m
\smallskip

\begin{center}
\needspace{6\baselineskip}
{\small
\begin{tabular}{L{0.17\linewidth} C{0.10\linewidth} C{0.10\linewidth} C{0.10\linewidth}
                C{0.10\linewidth} C{0.12\linewidth} C{0.12\linewidth}}
\toprule
\textbf{Fluid} & $t_1$, s & $t_2$, s & $t_3$, s & $t_4$, s & $t_\mathrm{avg}$, s & $v$, m/s \\
\midrule
1. & & & & & & \\ \rowrule
2. & & & & & & \\
\bottomrule
\end{tabular}\par
\vspace{2pt}
{\footnotesize\itshape Table 2. Fall times over the measured distance.}}
\end{center}

\begin{equation}
  t_\mathrm{avg} = \frac{t_1+t_2+t_3+t_4}{4}
  \qquad\qquad
  v = \frac{L}{t_\mathrm{avg}}
\end{equation}

\subsection{Viscosity and unit conversion}
Compute $\Delta\rho$, then apply Equation~(1). Convert the result to centipoise and to
centistokes. Show the substituted numbers, not only the answer.

\smallskip
\noindent Fluid 1: \quad $\Delta\rho =$ \bl\ $-$ \bl\ $=$ \bl\ kg/m$^3$ \par
\smallskip
\noindent\hspace{2em} $\mu =$ \bll\ Pa$\cdot$s $=$ \bll\ cP $=$ \bll\ cSt \par
\smallskip
\noindent Fluid 2: \quad $\Delta\rho =$ \bl\ $-$ \bl\ $=$ \bl\ kg/m$^3$ \par
\smallskip
\noindent\hspace{2em} $\mu =$ \bll\ Pa$\cdot$s $=$ \bll\ cP $=$ \bll\ cSt \par

\subsection{Uncertainty in the fall time}
Report the spread of your four readings, not only their average. The range divided by the
average, expressed as a percentage, is a defensible first estimate of the uncertainty in your
velocity, and therefore in your viscosity, since $\mu$ is inversely proportional to $v$.

\smallskip
\noindent Fluid 1: range $=$ \bls\ s \quad range$/t_\mathrm{avg} =$ \bls\ \% \quad$\rightarrow$\quad
$\mu =$ \bl\ $\pm$ \bl\ Pa$\cdot$s \par
\smallskip
\noindent Fluid 2: range $=$ \bls\ s \quad range$/t_\mathrm{avg} =$ \bls\ \% \quad$\rightarrow$\quad
$\mu =$ \bl\ $\pm$ \bl\ Pa$\cdot$s \par

\subsection{Is Stokes' law valid for your measurement?}
Equation~(1) is derived for \emph{creeping flow}, meaning the viscous forces on the sphere
overwhelm the inertial ones. The Reynolds number measures that ratio, and Stokes' law requires
it to be much less than one.

\begin{equation}
  \mathit{Re} = \frac{\rho_f\,v\,(2r)}{\mu}
\end{equation}

\noindent Fluid 1: $\mathit{Re} =$ \bl \quad Circle one:\quad
valid ($\mathit{Re}<1$) \quad$\cdot$\quad marginal (1 to 10) \quad$\cdot$\quad
invalid ($\mathit{Re}\geq10$) \par
\smallskip
\noindent Fluid 2: $\mathit{Re} =$ \bl \quad Circle one:\quad
valid ($\mathit{Re}<1$) \quad$\cdot$\quad marginal (1 to 10) \quad$\cdot$\quad
invalid ($\mathit{Re}\geq10$) \par

\smallskip
\noindent The Reynolds number above uses the viscosity you are testing, so when Stokes' law fails the
inflated $\mu$ makes $\mathit{Re}$ look smaller than it is. Recompute it with the rotary reading for
the same fluid at the same temperature (Part III). That is the honest check, and it is the value to
use in Questions~5 and~10.

\smallskip
\noindent Fluid 1: $\mathit{Re}$ with rotary $\mu$ $=$ \bl \qquad Fluid 2: $\mathit{Re}$ with rotary $\mu$ $=$ \bl \par

\begin{labnote}
\noteh{Do not discard a number because the model failed}
If your Reynolds number is above one, report the viscosity you calculated \emph{and} state that
Stokes' law does not strictly apply to it. A result reported with its limitation is worth more
than a result quietly deleted. Question~5 asks you to say which direction the error runs.
\end{labnote}

% =============================================================================
\needspace{8\baselineskip}
\section{Part II: Saybolt Viscometer}
The Saybolt viscometer measures fluid flow through a capillary of very small diameter. Before
filling the vessel, secure the rubber stopper at the base. The container must be filled just to
overflowing for accurate and consistent results, because the height of fluid above the orifice
sets the driving pressure. Position the beaker underneath and begin timing at the instant the
stopper is removed. Record the time in Saybolt Universal Seconds required for 60\,mL to collect
in the beaker.

Heat 60\,mL of each fluid to 50\,\textdegree C, pour it into the viscometer, and record the
collection time.

\begin{center}
\needspace{6\baselineskip}
{\small
\begin{tabular}{L{0.34\linewidth} C{0.26\linewidth} C{0.26\linewidth}}
\toprule
\textbf{Quantity} & \textbf{Fluid 1} & \textbf{Fluid 2} \\
\midrule
Fluid name                       & & \\ \rowrule
Temperature, \textdegree C       & & \\ \rowrule
Collection time, SUS             & & \\ \rowrule
Equation used, (6) or (7)        & & \\ \rowrule
$\nu$, cSt                       & & \\ \rowrule
$\rho$, g/cm$^3$                 & & \\ \rowrule
$\mu$, cP                        & & \\
\bottomrule
\end{tabular}\par
\vspace{2pt}
{\footnotesize\itshape Table 3. Saybolt measurements and conversions.}}
\end{center}

Convert the collection time to kinematic viscosity with the expression appropriate to your
reading, then to dynamic viscosity using the density you measured in Table~1.

\begin{align}
  \text{For } 32<\mathrm{SUS}<100:&\qquad \nu\,[\mathrm{cSt}] = 0.226\,\mathrm{SUS}-\frac{195}{\mathrm{SUS}}\\[2pt]
  \text{For } \mathrm{SUS}>100:&\qquad \nu\,[\mathrm{cSt}] = 0.220\,\mathrm{SUS}-\frac{135}{\mathrm{SUS}}
\end{align}
\begin{equation}
  \mu\,[\mathrm{cP}] = \nu\,[\mathrm{cSt}]\times\rho\,[\mathrm{g/cm^3}]
\end{equation}

\begin{labnote}
\noteh{The two expressions meet}
Both Equation~(6) and Equation~(7) return 20.65\,cSt at $\mathrm{SUS}=100$, so the pair is
continuous and the choice of branch at the boundary does not matter. Below 32\,s neither
applies; if your reading is that short, the fluid is too thin for this instrument and you should
say so rather than force a number out of Equation~(6).
\end{labnote}

% =============================================================================
\section{Part III: Rotary Viscometer}
A rotary viscometer comes closest to reporting the absolute viscosity of a fluid directly. A rotor immersed in the
sample is driven at a set speed, and the torque required to maintain that speed is proportional
to the fluid's viscosity. Sensors inside the instrument measure that torque and the viscosity is
displayed on the screen. The measurement is easier to perform and generally more accurate than
the other two methods, which is what makes it a useful reference against your falling-ball
result.

Identify the components of the instrument and become familiar with the values shown on the
display. Fill a beaker with 500\,mL of fluid and heat it to 60\,\textdegree C. Record the
temperature and the displayed viscosity in mPa$\cdot$s at 5\,\textdegree C intervals as the
sample cools, and convert each reading to Pa$\cdot$s using Equation~(9).

\begin{equation}
  \mu\,[\mathrm{Pa\cdot s}] = \frac{\mu\,[\mathrm{mPa\cdot s}]}{1000}
  \qquad (1\,\mathrm{Pa\cdot s} = 1000\,\mathrm{mPa\cdot s})
\end{equation}

Record the rotor and the speed before you take any reading. The torque the instrument senses
depends on the rotor diameter, its immersed height, the clearance to the beaker wall and the
rotation speed, so a viscosity value taken on this bench is only reproducible by someone else if
all of them are written down. A reading quoted without its spindle and speed is not a
measurement, it is a number.

\smallskip
\noindent Rotor diameter $D$: \bl\ mm \hfill Immersed height $H$: \bl\ mm \hfill
Speed: \bl\ rev/min \par
\smallskip
\noindent Beaker inner diameter: \bl\ mm \hfill Fluid depth: \bl\ mm \hfill
Annular gap $(B-D)/2$: \bl\ mm \par

\begin{center}
\needspace{6\baselineskip}
{\small
\begin{tabular}{C{0.09\linewidth} C{0.17\linewidth} C{0.20\linewidth}
                C{0.22\linewidth} C{0.20\linewidth}}
\toprule
\textbf{Trial} & \textbf{Set temp.} & \textbf{Measured temp.} &
\textbf{Viscosity} & \textbf{Viscosity} \\
 & \textdegree C & \textdegree C & mPa$\cdot$s & Pa$\cdot$s \\
\midrule
1 & 20 & & & \\ \rowrule
2 & 25 & & & \\ \rowrule
3 & 30 & & & \\ \rowrule
4 & 35 & & & \\ \rowrule
5 & 40 & & & \\ \rowrule
6 & 45 & & & \\ \rowrule
7 & 50 & & & \\ \rowrule
8 & 55 & & & \\ \rowrule
9 & 60 & & & \\
\bottomrule
\end{tabular}\par
\vspace{2pt}
{\footnotesize\itshape Table 4. Rotary viscometer readings. Fluid type: \blll}}
\end{center}

\subsection{What the instrument computes}
The display gives you a viscosity directly, but it is worth knowing what stands behind that
number, because it explains why the spindle and the speed must be recorded. Treat the rotor as a
cylinder of radius $R_i$ turning at angular velocity $\omega$ inside a stationary beaker of
inner radius $R_o$, immersed to a depth $H$. The instrument holds the speed constant and
measures the torque $M$ required to maintain it.

\begin{equation}
  \omega = \frac{2\pi N}{60}
  \qquad\qquad
  \dot{\gamma} = \frac{2\,\omega\,R_o^{2}}{R_o^{2}-R_i^{2}}
\end{equation}
\begin{equation}
  \tau = \frac{M}{2\pi R_i^{2} H}
  \qquad\qquad
  \mu = \frac{\tau}{\dot{\gamma}}
\end{equation}

Combining the two gives the viscosity in terms of the measured torque alone:

\begin{equation}
  \mu = \frac{M\,(R_o^{2}-R_i^{2})}{4\pi\,\omega\,H\,R_i^{2}R_o^{2}}
\end{equation}
\noindent{\small
\begin{tabular}{@{\hspace{1.5em}}l@{\quad}l@{\hspace{2em}}l@{\quad}l@{}}
$N$          & rotation speed, rev/min      & $\omega$ & angular velocity, rad/s \\[-6pt]
$R_i$        & rotor radius, m              & $R_o$    & beaker inner radius, m \\[-6pt]
$H$          & immersed height, m           & $M$      & torque on the rotor, N$\cdot$m \\[-6pt]
$\dot\gamma$ & shear rate, s$^{-1}$         & $\tau$   & shear stress, Pa \\
\end{tabular}\par}

\smallskip
Every quantity in Equation~(12) except $M$ and $N$ is fixed once the spindle and beaker are
chosen, so the whole expression collapses to $\mu = KM/N$, where $K$ is a pure geometric
constant. That constant is the spindle factor printed in the instrument's tables, and it is the
reason a reading is meaningless unless the spindle and the speed are recorded alongside it.

\begin{labnote}
\noteh{Assumptions behind Equation~(12)}
The fluid is Newtonian, the flow is laminar, the rotor is fully immersed and turning coaxially
in the beaker, and the torque contributed by the top and bottom faces of the rotor is neglected.
Equation~(12) is exact for a long cylinder in a close-fitting beaker and becomes an
approximation as the gap widens.
\end{labnote}

\smallskip
\noindent Compute for one of your readings: \quad
$\omega =$ \bl\ rad/s \quad $\dot\gamma =$ \bl\ s$^{-1}$ \quad $M =$ \bl\ N$\cdot$m \par

% =============================================================================
\needspace{14\baselineskip}
\subsection{Temperature plot}
Plot temperature on the horizontal axis and viscosity on the vertical axis. Sketch the curve
here during the laboratory period, then produce the same plot in a spreadsheet and submit it
with your report. Label both axes with units.

\begin{center}
  \graphpaper{16}{9.5}
\end{center}
\vspace{-4pt}
\noindent{\footnotesize\color{ink2}
20\,\textdegree C \hfill Temperature \hfill 60\,\textdegree C\par}

% =============================================================================
\section{Analysis Questions}
\begin{enumerate}
  \item Which of your two fluids has the greater resistance to flow? Cite the specific measured
        quantity that establishes this, not the calculated viscosity alone.
        \answerline\answerline

  \item Your falling-ball result is a dynamic viscosity and your Saybolt result begins as a
        kinematic viscosity. Explain why the Saybolt method yields kinematic viscosity directly,
        given that the fluid is driven through the orifice by its own weight.
        \answerline\answerline\answerline

  \item Compare the viscosity you obtained for the same fluid by the falling-ball method and by
        the rotary viscometer at the same temperature. Express the disagreement as a percentage
        of the rotary value.
        \answerline\answerline

  \item List two experimental causes of the disagreement in Question~3 that are independent of
        the validity of Stokes' law.
        \answerline\answerline\answerline

  \item Using your Reynolds numbers from Section~5.6, state whether Stokes' law was applicable
        to each fluid. For any fluid where it was not, state the \emph{direction} of the
        resulting error: at Reynolds numbers above one, drag on the sphere is larger than
        Equation~(1) assumes, so does your computed viscosity read high or low? Justify from
        Equation~(1).
        \answerline\answerline\answerline

  \item Propose one change to the falling-ball apparatus that would bring the Reynolds number
        into the valid range, and estimate the factor by which it would reduce $\mathit{Re}$.
        Note that $v$ is proportional to $r^2$, so $\mathit{Re}$ is proportional to $r^3$.
        \answerline\answerline

  \item Describe the shape of your temperature-viscosity curve. Is a straight line an adequate
        description of it across the whole range? What does your answer imply about quoting a
        single viscosity value for a hydraulic fluid without stating a temperature?
        \answerline\answerline\answerline

  \item A hydraulic system designed around a fluid at 40\,\textdegree C is started on a winter
        morning at 5\,\textdegree C. Using the trend in your data, argue whether the pump inlet
        or the actuator is the more likely place for trouble, and say what physically goes wrong
        there.
        \answerline\answerline\answerline

  \item An engine oil is labelled 15W-40 and a hydraulic fluid is labelled ISO~46. State which
        viscosity family each grading system refers to and at what temperature, and explain why
        the two labels cannot be compared directly.
        \answerline\answerline\answerline

  \item Of your three measurements of the same fluid, which do you trust most? Defend the choice
        using your uncertainty estimate from Section~5.5 and your Reynolds numbers from
        Section~5.6, not the reputation of the instrument.
        \answerline\answerline\answerline
\end{enumerate}

% =============================================================================
\section{Interactive Simulator}
An interactive version of all three instruments is available in the browser. It runs the same
equations printed in this handout, animates the sphere falling at the velocity your own timings
imply, drains the Saybolt cup over your own SUS reading, and plots your rotary data as you enter
it. It computes the Reynolds number live and labels the Stokes regime, so you can see directly
how sphere size moves a measurement in and out of validity. Use it to check your arithmetic and
to explore fluids the bench does not stock. It does not replace the bench measurements.

\smallskip
\noindent\small\url{https://alibulentkoc.github.io/viscosity-measurement-lab/lab/viscosity-lab.html}

\vfill
\noindent\rule{\linewidth}{0.4pt}
\noindent{\footnotesize\color{ink2}
\textbf{Cite this activity as:} Koc, A.~B. \textit{Viscosity Measurement Laboratory: Falling
Ball, Saybolt, and Rotary Viscometers}. Department of Agricultural Sciences, Agricultural
Mechanization and Business Program, Clemson University.
DOI \texttt{10.5281/zenodo.22218165} \quad$\cdot$\quad ORCID \texttt{0000-0003-3633-1699}.\par
\smallskip
Laboratory materials licensed CC BY 4.0; simulator source licensed MIT.\par}

\end{document}
